\subsection{From block libraries to image unions above dimension
48}

The initial family used 63 disjoint blocks of a Pless neighbour. Its
prefix sizes, combined with root tail codes, produced seven
configurations in dimensions 49 through 55. The paper retains the
neighbour's construction and minimum-norm proof.

The larger family uses the published catalogue basis of \(P_{48p}\). A
public class of 7,069 minimal lines was extended to 7,077. Its
automorphism images were then selected jointly, using the size of their
union as the objective. Assigning repeated lines to their first
occurrence produced disjoint subsets without losing the within-class
product bound.

The seven union sizes are

\[
7077,\ 21231,\ 42459,\ 84874,\ 141328,\ 253851,\ 442507.
\]

Each contributes two points per line, in addition to the mother shell
and the root tail code. The dimension-50 witness uses three disjoint
images. Larger image families retain some overlaps, whose exact effect
is recorded in the union sizes.

Different assigned classes use distinct antipodal tail lines. This gives
the cross-class product bound \((3/4)(1/2)+(1/4)(1/2)=1/2\). An
injective tail assignment therefore accompanies the union of the
selected head classes.
